%=====================================================================
\chapter{Transactions, Concurrency and Recovery}\label{ch:transactions}
%=====================================================================
\locator{6}{Part VI --- Making It Work, and What Comes Next}

\begin{outcomes}
By the end of this chapter you will be able to:
\begin{tight}
  \item \textbf{State} the four ACID properties and \textbf{give} a failure
        that each one prevents.
  \item \textbf{Show} a lost update happening, and \textbf{prevent} it with
        the right kind of transaction.
  \item \textbf{Explain} SQLite's locks and why it allows one writer at a time.
  \item \textbf{Distinguish} the deadlock SQLite can have from the one a server
        DBMS has, and \textbf{say} how each is resolved.
  \item \textbf{Back up} a live database safely, and \textbf{state} what the
        journal does after a crash.
\end{tight}
\end{outcomes}

\begin{prereq}
\chref{whynotfiles}: atomicity and concurrent access were two of the six
problems. \chref{whodoeswhat}: the transaction, lock and recovery managers.
\chref{evolution}: \sqlkw{BEGIN} and \sqlkw{COMMIT} in a migration. No
operating-systems course is assumed. You need
\texttt{code/m22\_transactions.py}.
\end{prereq}

\begin{keyterms}
transaction $\bullet$ \sqlkw{BEGIN} $\bullet$ \sqlkw{COMMIT} $\bullet$
\sqlkw{ROLLBACK} $\bullet$ \sqlkw{SAVEPOINT} $\bullet$ ACID $\bullet$
atomicity $\bullet$ consistency $\bullet$ isolation $\bullet$ durability
$\bullet$ autocommit $\bullet$ lost update $\bullet$ dirty read $\bullet$
serialisable $\bullet$ lock $\bullet$ \sqlkw{SHARED} $\bullet$
\sqlkw{RESERVED} $\bullet$ \sqlkw{EXCLUSIVE} $\bullet$ \sqlkw{DEFERRED}
$\bullet$ \sqlkw{IMMEDIATE} $\bullet$ deadlock $\bullet$ \texttt{SQLITE\_BUSY}
$\bullet$ \texttt{busy\_timeout} $\bullet$ rollback journal $\bullet$ WAL
$\bullet$ \texttt{synchronous}
\end{keyterms}

\section{Promises About Groups of Changes}

Every constraint in Part~II was a promise about a row. A transaction is a
promise about a \emph{group} of changes: that they happen together, or not at
all, and that nothing anybody else does at the same moment can see them half
done. \chref{whynotfiles} listed atomicity and concurrent access among the six
problems of files. This chapter is where the database keeps those promises ---
and measures what keeping them costs.

\begin{definitionbox}[Transaction, and ACID]
A \term{transaction} is a sequence of operations treated as one unit, begun
with \sqlkw{BEGIN} and ended with \sqlkw{COMMIT} (keep it) or \sqlkw{ROLLBACK}
(undo it all). Outside one, SQLite is in \term{autocommit}: every statement is
its own transaction.

\smallskip
\textbf{Atomicity}: all of it happens, or none. \textbf{Consistency}: it takes
the database from one state satisfying every constraint to another.
\textbf{Isolation}: concurrent transactions do not see each other's
unfinished work. \textbf{Durability}: once committed, it survives a crash.
\end{definitionbox}

\begin{measured}[A crash, a lost update, a deadlock, and the price of durability]
\textsf{\footnotesize \texttt{code/m22\_transactions.py}; two real connections
to one file, real failures. SQLite 3.51.2, Intel i5-6600, Windows 11,
mechanical disk.}

\medskip
\textbf{1. Atomicity} --- move a student between sections, a \sqlkw{DELETE}
then an \sqlkw{INSERT}, and fail between them

\begin{tabular}{@{}L{4.4cm}L{3.4cm}L{6.0cm}@{}}
& \textbf{enrolled in} & \\
autocommit & \textbf{nothing} & the \sqlkw{DELETE} was already permanent \\
\sqlkw{BEGIN} \dots\ \sqlkw{ROLLBACK} & section 1 & the two were one unit \\
\end{tabular}

\medskip
\textbf{2. The lost update} --- two clerks each add one to a seat count of 40:
A reads, B reads, A writes, B writes

\begin{tabular}{@{}L{3.4cm}L{8.6cm}R{1.6cm}@{}}
& \textbf{what happened} & \textbf{final} \\
no transaction & both read 40, both wrote 41 & \textbf{41} \\
\sqlkw{BEGIN DEFERRED} & B's write refused: \texttt{database is locked}; then
  A's \sqlkw{COMMIT} refused too; B rolled back; A committed & 41 \\
\sqlkw{BEGIN IMMEDIATE} & B could not begin until A had finished; then read 41
  and wrote 42 & \textbf{42} \\
\end{tabular}

\smallskip
With no transaction an enrolment simply vanished, and nobody was told. With
\sqlkw{DEFERRED} the two clerks \textbf{deadlocked}: B held a read lock and
wanted to write; A held the write lock and needed B's read lock released before
it could commit. Neither could move until B gave up --- so B's +1 was lost
again, \emph{but B was told}, and can retry. \sqlkw{IMMEDIATE} counted both.

\medskip
\textbf{3. \texttt{SQLITE\_BUSY}} --- A holds the write lock for 300\,ms; B
tries to write

\begin{tabular}{@{}L{4.4cm}L{6.0cm}R{3.0cm}@{}}
\textbf{B's busy\_timeout} & \textbf{outcome} & \textbf{B waited} \\
0\,ms (the default) & failed at once: \texttt{database is locked} & 0\,ms \\
1{,}000\,ms & succeeded after A committed & 455\,ms \\
\end{tabular}

\smallskip
B waited for A's 300\,ms, A's commit, and the interval at which SQLite's busy
handler retries --- and then went ahead.

\medskip
\textbf{4. Durability, priced} --- 200 inserts

\begin{tabular}{@{}L{7.4cm}R{2.8cm}R{2.8cm}@{}}
& \textbf{total} & \textbf{per row} \\
a commit each, \texttt{synchronous = FULL} & 4{,}590.9\,ms & 22.954\,ms \\
one commit for all, \texttt{synchronous = FULL} & 27.8\,ms & 0.139\,ms \\
a commit each, \texttt{synchronous = OFF} & 239.2\,ms & 1.196\,ms \\
\end{tabular}

\smallskip
\textbf{One commit for all 200 was 165$\times$ faster than one each.} A commit
waits until the disk confirms the data is really on it; on a mechanical disk
that wait is 23\,ms and is almost the whole cost. \texttt{synchronous = OFF}
skips the wait --- 19$\times$ faster, still 9$\times$ slower than batching ---
and gives up durability: a power cut can lose transactions that were reported
as committed, or corrupt the file. Batching keeps the guarantee \emph{and} is
faster.
\end{measured}

\begin{alertbox}[Four rules from the measurement]
\begin{tight}
  \item \textbf{Anything that is two statements and must be one fact goes in a
        transaction.} Measured: autocommit left a student enrolled in nothing.
  \item \textbf{A read-then-write goes in \sqlkw{BEGIN IMMEDIATE}.} It takes
        the write lock before the read, so nobody else can read the old value
        in between. \sqlkw{DEFERRED} (the default) deadlocked.
  \item \textbf{Set \texttt{busy\_timeout} on every connection.} The default
        is 0, so a second writer is refused instead of waiting. It is per
        connection, like \chref{constraints}'s \texttt{foreign\_keys}, and for
        the same reason has to be set every time.
  \item \textbf{Group writes into transactions; never turn durability off to
        go faster.} The guarantee and the speed are not in conflict here.
\end{tight}
\end{alertbox}

\section{How SQLite Isolates: One Lock for the Whole File}

\begin{definitionbox}[SQLite's locks]
SQLite locks the \emph{database file}, not rows or tables. A connection
progresses through:

\smallskip
\begin{tabular}{@{}L{2.8cm}L{11.0cm}@{}}
\sqlkw{SHARED} & may read. Any number of connections may hold it. \\
\sqlkw{RESERVED} & intends to write. Only one at a time; readers continue. \\
\sqlkw{EXCLUSIVE} & is committing. Needed to write to the file; requires
  every \sqlkw{SHARED} lock to be released first. \\
\end{tabular}

\smallskip
\sqlkw{BEGIN DEFERRED} takes no lock until the first read (then
\sqlkw{SHARED}) or write (then \sqlkw{RESERVED}). \sqlkw{BEGIN IMMEDIATE}
takes \sqlkw{RESERVED} at once. \sqlkw{BEGIN EXCLUSIVE} takes
\sqlkw{EXCLUSIVE} at once.

\smallskip
The consequence: \textbf{one writer at a time, for the whole database}, and
transactions that are \term{serialisable} --- the result is always as if they
had run one after another. There are no weaker isolation levels to choose
from, because there is nothing to relax --- and so no \term{dirty read}:
no connection can ever see another's uncommitted change.

\smallskip
A \term{savepoint} marks a point inside a transaction to roll back to
without abandoning the whole of it: \texttt{SAVEPOINT s1;} \dots\
\texttt{ROLLBACK TO s1;} undoes only the work since \texttt{s1}.
\end{definitionbox}

\dsfig{%
\begin{tikzpicture}[
  tx/.style={rounded corners=3pt, draw=#1, line width=1pt, fill=#1!8,
             text=#1!60!black, font=\scriptsize\bfseries, align=center,
             minimum width=2.2cm, minimum height=0.8cm},
  rs/.style={rounded corners=3pt, draw=neutral, line width=0.9pt,
             fill=neutral!8, text=neutral!80!black, font=\scriptsize,
             align=center, minimum width=2.4cm, minimum height=0.8cm},
  hold/.style={-{Stealth[length=2.2mm]}, greenhl!70!black, line width=1.1pt},
  want/.style={-{Stealth[length=2.2mm]}, accent, line width=1.1pt, dashed},
  ttl/.style={font=\scriptsize\bfseries, text=primary, anchor=south},
  lb/.style={font=\scriptsize\itshape, text=#1}
]
  % --- SQLite: one file, a lock upgrade
  \node[ttl] at (-3.9,1.85) {SQLite: one lock, upgraded};
  \node[tx=secondary] (a) at (-5.6,0.9) {A};
  \node[tx=purplehl]  (b) at (-2.2,0.9) {B};
  \node[rs] (f) at (-3.9,-0.6) {the database file};
  \draw[hold] (a.south) -- node[lb=greenhl!60!black, left, pos=0.45]
       {holds RESERVED} (f.north west);
  \draw[hold] (b.south) -- node[lb=greenhl!60!black, right, pos=0.45]
       {holds SHARED} (f.north east);
  \draw[want] (a.east) -- node[lb=accent, above] {needs B to let go} (b.west);
  \node[lb=neutral, text width=6.0cm, align=center] at (-3.9,-1.55)
       {A cannot commit while B reads; B cannot write while A holds
        RESERVED. \textbf{Measured} --- SQLite refuses, one rolls back.};

  % --- server: two rows, a cycle
  \node[ttl] at (3.9,1.85) {a server DBMS: two rows, a cycle};
  \node[tx=secondary] (c) at (2.2,0.9) {T1};
  \node[tx=purplehl]  (d) at (5.6,0.9) {T2};
  \node[rs] (r1) at (2.2,-0.6) {row: section 1};
  \node[rs] (r2) at (5.6,-0.6) {row: section 2};
  \draw[hold] (c) -- node[lb=greenhl!60!black, left] {locks} (r1);
  \draw[hold] (d) -- node[lb=greenhl!60!black, right] {locks} (r2);
  \draw[want] (c.south east) -- node[lb=accent, above, sloped, pos=0.35]
       {wants} (r2.north west);
  \draw[want] (d.south west) -- node[lb=accent, above, sloped, pos=0.35]
       {wants} (r1.north east);
  \node[lb=neutral, text width=6.0cm, align=center] at (3.9,-1.55)
       {Each waits for the other's row for ever. The server finds the cycle
        and kills one. \textbf{Cannot happen in SQLite}: there are no row
        locks.};
\end{tikzpicture}}%
{Two kinds of deadlock. The left one was produced by the harness with two
\sqlkw{DEFERRED} transactions. The right one needs row-level locks, which
SQLite does not have, so it is drawn here rather than run --- \chref{project}
is where a server's locking is compared properly.}{deadlock}

\begin{conceptbox}[What SQLite does not have, and what you would get instead]
A server DBMS --- PostgreSQL, MySQL, Oracle --- locks rows, and most use
\term{multi-version concurrency control}: a writer creates a new version of a
row while readers keep seeing the old one, so readers and writers never block
each other. It offers several \emph{isolation levels}, from \textsc{read
committed} to \textsc{serializable}, trading correctness for concurrency.

\smallskip
SQLite has none of that, and for one database file on one machine it does not
need it. What it costs is concurrency: many readers, \textbf{one} writer. For a
university's results system with a few clerks it is plenty; for a booking site
with a thousand simultaneous purchases it is the reason to move to a server,
and \chref{project} makes that comparison.
\end{conceptbox}

\section{Recovery, and Backup}

\begin{definitionbox}[The journal, and what happens after a crash]
Before changing a page of the database file, SQLite writes the page's
\emph{original} contents to a \term{rollback journal}, \texttt{university.db-journal}.
Only then does it change the file. At \sqlkw{COMMIT} it deletes the journal.

\smallskip
After a crash, the next connection finds a journal still present, knows a
transaction was interrupted, and copies the original pages back. The database
returns to its state before the transaction began --- atomicity and
durability, kept by one extra file.

\smallskip
\term{WAL} mode (\texttt{PRAGMA journal\_mode = WAL}) reverses the idea:
changes are appended to a write-ahead log and copied into the database later.
Readers then never block the writer, which is the nearest SQLite comes to a
server's concurrency.
\end{definitionbox}

\begin{examplebox}[Backing up a database that is in use]
\textbf{The wrong way:} \texttt{copy university.db backup.db} while a program
is writing. If the copy happens mid-transaction it captures a half-written
file without its journal --- a backup that may not open, discovered on the day
it is needed.

\smallskip
\textbf{Two right ways}, both of which take a consistent snapshot:
\begin{lstlisting}[language=sql]
.backup backup.db                     -- in the sqlite3 shell
VACUUM INTO 'backup.db';              -- as SQL, from any program
\end{lstlisting}

\smallskip
And \chref{evolution}'s rule applies: \textbf{open the backup} and run
\texttt{PRAGMA integrity\_check;} on it. A backup that has never been restored
is a hope.
\end{examplebox}

\begin{pitfall}
\begin{tight}
  \item \textbf{Two statements that must be one fact, in autocommit.} Measured:
        a student enrolled in nothing.
  \item \textbf{Read-then-write in \sqlkw{BEGIN DEFERRED}.} Measured: a
        deadlock. Use \sqlkw{IMMEDIATE}.
  \item \textbf{Leaving \texttt{busy\_timeout} at 0.} A second writer is
        refused instead of waiting.
  \item \textbf{A commit per row in a bulk load.} Measured: 165$\times$ slower.
  \item \textbf{\texttt{synchronous = OFF} for speed.} It gives up durability
        and is still slower than batching.
  \item \textbf{Copying a live database file as a backup.} Use
        \texttt{.backup} or \sqlkw{VACUUM INTO}, then open the result.
  \item \textbf{Python's default \texttt{isolation\_level}.} The
        \texttt{sqlite3} module silently issues \sqlkw{BEGIN} for you before
        writes. The harness sets \texttt{isolation\_level=None} so every
        \sqlkw{BEGIN} is one you wrote.
\end{tight}
\end{pitfall}

\begin{lab}[Lose an enrolment, deadlock two clerks, and back up safely]
\begin{lstlisting}[language=bash]
# --- 1. Two shells on ONE file: open two terminals ------------------
sqlite3 txn.db          # terminal A
sqlite3 txn.db          # terminal B
\end{lstlisting}
\begin{lstlisting}[language=sql]
-- --- 2. In A only: a table to fight over ------------------------
CREATE TABLE section (id TEXT PRIMARY KEY, seats_taken INTEGER) STRICT;
INSERT INTO section VALUES ('CS-103/1', 40);

-- --- 3. The lost update, by hand: type in this order -------------
-- A:  SELECT seats_taken FROM section;         -- 40
-- B:  SELECT seats_taken FROM section;         -- 40
-- A:  UPDATE section SET seats_taken = 41;
-- B:  UPDATE section SET seats_taken = 41;
-- A:  SELECT seats_taken FROM section;         -- 41. Two clerks, +1.

-- --- 4. The deadlock: reset to 40, then type in this order -------
-- A:  BEGIN DEFERRED;  SELECT seats_taken FROM section;
-- B:  BEGIN DEFERRED;  SELECT seats_taken FROM section;
-- A:  UPDATE section SET seats_taken = 41;     -- works
-- B:  UPDATE section SET seats_taken = 41;     -- database is locked
-- A:  COMMIT;                                  -- also locked. Why?
-- B:  ROLLBACK;
-- A:  COMMIT;                                  -- now it works

-- --- 5. The fix: reset to 40, then ------------------------------
-- A:  BEGIN IMMEDIATE;  SELECT seats_taken FROM section;
-- B:  BEGIN IMMEDIATE;                         -- refused: wait
-- A:  UPDATE section SET seats_taken = 41;  COMMIT;
-- B:  BEGIN IMMEDIATE;  SELECT seats_taken FROM section;   -- 41
-- B:  UPDATE section SET seats_taken = 42;  COMMIT;

-- --- 6. Waiting instead of failing ------------------------------
-- B:  .timeout 5000
-- Repeat step 5. What does B do now instead of refusing?

-- --- 7. A safe backup -------------------------------------------
.backup txn_backup.db
\end{lstlisting}
\begin{lstlisting}[language=bash]
# --- 8. Open the backup and check it --------------------------------
sqlite3 txn_backup.db "PRAGMA integrity_check; SELECT * FROM section;"

# --- 9. The measured version ----------------------------------------
python code/m22_transactions.py
\end{lstlisting}

\smallskip
\textbf{What to notice.} Step 4's second \sqlkw{COMMIT} is the one to explain.
A already holds the write lock; to commit it needs the \sqlkw{EXCLUSIVE} lock,
which requires every reader to leave --- and B is a reader, waiting to become a
writer. Neither can move. SQLite cannot wait its way out of that, so it
refuses, and one of the two must roll back.
\end{lab}

\begin{checkpoint}
\begin{tightnum}
  \item For each ACID property, give the measured result in this chapter that
        shows it being kept or broken, and say what would have happened
        without it.
  \item Two clerks deadlocked under \sqlkw{BEGIN DEFERRED} and both increments
        counted under \sqlkw{BEGIN IMMEDIATE}. Explain both outcomes in terms of
        which lock each connection held and wanted at each step.
  \item A bulk load of 200 rows took 4.6 seconds with a commit per row and
        28\,ms in one transaction. Explain the difference, say why
        \texttt{synchronous = OFF} is not the fix, and say what it would cost
        the university if it were used.
\end{tightnum}
\end{checkpoint}

\begin{chaptersummary}
\begin{tight}
  \item A \textbf{transaction} groups changes into one unit with
        \sqlkw{BEGIN}, \sqlkw{COMMIT} and \sqlkw{ROLLBACK}. Outside one, every
        statement is its own (\textbf{autocommit}).
  \item \textbf{ACID}: atomicity, consistency, isolation, durability.
  \item \textbf{Measured:} a failure between a \sqlkw{DELETE} and an
        \sqlkw{INSERT} left a student \textbf{enrolled in nothing} in
        autocommit; inside a transaction, nothing changed.
  \item \textbf{Measured:} two read-then-write clerks \textbf{lost an update}
        with no transaction, \textbf{deadlocked} under \sqlkw{DEFERRED}, and
        were both counted under \sqlkw{IMMEDIATE}.
  \item SQLite locks the whole file --- \sqlkw{SHARED}, \sqlkw{RESERVED},
        \sqlkw{EXCLUSIVE} --- so it has one writer at a time and is always
        serialisable.
  \item \textbf{Measured:} with \texttt{busy\_timeout} 0 a second writer was
        refused; with 1{,}000\,ms it waited 455\,ms and succeeded. Set it on
        every connection.
  \item A server's two-row deadlock cycle cannot occur in SQLite, which has no
        row locks; SQLite's deadlock is two connections on the one file lock.
  \item \textbf{Measured:} one commit for 200 inserts was \textbf{165$\times$}
        faster than one each. \texttt{synchronous = OFF} gave up durability and
        was still 9$\times$ slower than batching.
  \item The \textbf{rollback journal} holds original pages until commit, and is
        replayed after a crash. Back up with \texttt{.backup} or
        \sqlkw{VACUUM INTO}, never by copying a live file --- then open the
        backup.
\end{tight}
\end{chaptersummary}

\begin{reviewq}
  \item \textbf{(MCQ)} In autocommit mode, each statement is: (a)~never
        committed; (b)~its own transaction; (c)~part of one transaction per
        connection; (d)~rolled back on exit.
  \item \textbf{(MCQ)} \sqlkw{BEGIN IMMEDIATE} acquires: (a)~no lock;
        (b)~\sqlkw{SHARED}; (c)~\sqlkw{RESERVED}; (d)~a row lock.
  \item \textbf{(Short)} Why does a read-then-write need its write lock before
        the read?
  \item \textbf{(Short)} Why can SQLite never have a two-row deadlock cycle?
  \item \textbf{(Short)} What does the rollback journal contain, and what is
        done with it after a crash?
  \item \textbf{(Applied)} Write the transaction that moves a student from one
        section to another and adjusts both sections' seat counts. Say which
        kind of \sqlkw{BEGIN} you would use and why, what happens if it fails
        after the second statement, and how a second clerk doing the same move
        at the same moment would be affected.
  \item \textbf{(Applied)} The university's results are entered by six clerks
        at once during the last week of term. Using the measured lock
        behaviour, explain what each clerk would experience with
        \texttt{busy\_timeout} 0 and with 5{,}000\,ms, estimate how long a
        clerk might wait if each commit took 23\,ms, and say at what number of
        clerks you would recommend moving to a server DBMS --- and why the
        answer is not a fixed number.
\end{reviewq}
